\scanpage{10}
% archive sheet number: 9
% lines 1–4
déjà que le foncteur $M\unsure{\leadsto|\sim|\mapsto}\widetilde M$ des $\Gamma B$-\deleted{\unsure{M|$\mathcal M$}}modules de présentation finie dans les $B$-Modules de prés. finie est pl. fidèle.

% lines 4–8
\deleted{\unsure{Pour voir|Prouvons} Reste à prouver que si $M$ est un $B$-Module de prés. finie, il est engendré par ses sections. Soit $x$ un pt fermé de $X$ [\inserted{N.B.} (Nous supposons $X$ \textquotedblleft pur\textquotedblright{} ! ]. Alors $M_x$ est un $B_x$-Module de présentation finie}

% lines 9
\deleted{Soit $M'$ le sous-module engendré par les sections.}

% lines 10–14
(en utilisant \unsure{seult} la condition $\alpha$) de la \textquotedblleft remarque\textquotedblright{} ci-dessus). Reste à prouver $\beta$) [qui \unsure{en conclut|ne cloche|est inclus} dans le cas $\underline O_X$ \unsure{fixé|fini} sur \unsure{$X$|$\hat x$} compact \ldots !.], sous les conditions plus fortes \unsure{précisées|possibles}.

% lines 14–17
Il faut prouver que $B$ satisfait conditions (ii) et (iii), pour (ii) c'est l'hyp, puis (iii)$'$ c'est encore un passage à la limite, car

% lines 18
\[
 J=\varinjlim_m J_m
\]

% lines 19
il faut prouver que les $J_m\inserted{=J\cap B_m}$ sont cohérents
% slanted
\margin{? histoire d'anneaux cohérents ??}

% lines 20
-----
